\clearpage
\fancyhead[L]{\small HERSCHEL AND SOUTH. OBSERVATIONS OF DOUBLE STARS.}
\hypertarget{native-647-lower}{}
\begin{center}\small Göttingische gelehrte Anzeigen. No. 201. pp. 2001..2004. 18 December 1826.\\[3pt]\rule{.72\linewidth}{.3pt}\end{center}

\emph{Observations of the apparent distances and positions of \textup{380} double and triple stars made in the years \textup{1821}, \textup{1822} and \textup{1823}, and compared with those of other astronomers.} London 1825. By \textsc{J. F. W. Herschel} and \textsc{James South}. 412 pages in quarto, with 4 copperplates.

\emph{Observations of the apparent distances and positions of \textup{458} double and triple stars made in the years \textup{1823}, \textup{1824} and \textup{1825}.} Ibid. 1826. By \textsc{James South}. 391 and XVIII pages in quarto.

The history of astronomy often shows us that important and unexpected insights were gained from observations that had been made and recorded with painstaking care long before, without any inkling of such an outcome. Double stars are a striking example. A number of double stars, albeit a comparatively small one, were already known before \textsc{Herschel}, but they received no special attention. The \SourceLeaf{648}view was that the phenomenon of a double star was merely a play of chance, which had placed two fixed stars in nearly the same direction from our solar system, while their actual distances might in general differ greatly. \textsc{Herschel} was the first to call attention to the advantages that continued observation of a double star might offer for determining the annual parallax, a method whose usefulness plainly depends entirely on that assumption. In his survey of the starry sky, \textsc{Herschel} recorded every double star he encountered, observed the relative positions with great accuracy, and thus compiled a catalogue of more than 700 double and multiple stars. Many of these, to be sure, can be called double stars only loosely, owing to the considerable separation of their component stars, though the nature of the matter precludes drawing any sharp limits with respect to that separation.

Twenty years later, \textsc{Herschel} re-examined the double stars. In many of them he noticed changes in the relative positions of the individual stars that led to the conviction that the earlier view of the chance formation of double stars could not hold generally, and that at least many of them might be pairs of stars \emph{really} close together, moving around one another according to definite laws grounded in the forces of nature.

This highly important conclusion was made much more probable by the very large number of double stars observed in the sky. Yet one may boldly assert that all this was hardly needed to expose the weakness of the earlier assumption. Indeed, consider for example just the one double star Castor. It consists of two stars that differ little in brightness, are both about third magnitude, and stand about five seconds of arc apart. If we assume an entirely random distribution on the celestial sphere of all fixed stars down to the third magnitude, we find the probability of such a double star forming to be so extraordinarily small that, in ordinary life, in analogous circumstances, we would consider it absurd to doubt that its origin was \emph{not} due to chance. It is of course possible that among the great number of double stars there may be some that owe their origin to chance. But we must suppose that in most of the closer pairs the stars really belong to the same system, and that their relative motions have opened a new field for interesting observations. These motions are so slow, however, that for most double stars the displacement is slight even over a generation. Yet some double stars have been found for which it becomes visible in comparatively very short intervals. In any event, it is clear that re-examining the double stars from time to time is of great importance if we are to prepare material for future centuries to gain a vastly greater knowledge of the structure of the universe.

Messrs. \textsc{Herschel} the younger, son of the famous observer, and \textsc{South} have done this in an exemplary manner in the works under review. The first work contains the observations these two astronomers made jointly with a five-foot and a seven-foot equatorial telescope. The second contains those the latter astronomer made alone, partly with both instruments in England and partly with the seven-foot equatorial telescope in France, at an observatory temporarily set up in Passy. We thus have a wealth of observations of 838 double stars from the years 1821..1825, together with their comparison with the earlier observations of the elder \textsc{Herschel}, partly corrected and supplemented from the manuscripts he left behind, and with the no less valuable observations of \textsc{Struve}. For most double stars the changes in position observed so far are small, though \SourceLeaf{649}unmistakable. The period of one sun's orbit around the other is generally measured in centuries. Yet in some cases the changes are already large enough to give us an approximate idea of the length of the period. For the remarkable star 61 in the Swan, for example, the change in position angle from 1753 to 1825 already amounted to 53 degrees. For Castor, from 1759 to the same year, the change was 63 degrees. For $\sigma$ in the Crown since 1781 it was 90 degrees, for $\xi$ in the Great Bear over that same interval it was 259 degrees, and for 70 Ophiuchus since 1779 it was as much as 302 degrees. Our limited space prevents us from pointing out many other noteworthy cases. There are stars that once appeared double and now coincide completely, triple stars that form One system and move relative to one another, stars that stand twelve minutes of arc apart but very probably form only One system, and the like. [See Neue Aussicht zur Erweiterung etc. \hyperlink{native-186}{above, p. 181.} \textsc{Schering}.]
